\documentclass[11pt, a4paper]{article}

\input{bfw-style.tex}

\title{\vspace*{-1.5cm}\Huge\bfseries\color{bfwPrimaryDark}Aufgabenheft Session~8\\[0.4em]
  \large\color{bfwInkSoft}\textnormal{Präsentieren, Feedback \& Gesamtwiederholung --- Übungen im IHK-Klausurformat}}
\author{\textsf{IT Lernfeld 1 --- Das Unternehmen und die eigene Rolle im Betrieb beschreiben}}
\date{Selbstlernmaterial}

\begin{document}
\maketitle
\thispagestyle{empty}

\begin{merksatz}[title={\bfseries So nutzen Sie dieses Aufgabenheft}]
Bearbeiten Sie alle Aufgaben zunächst \emph{ohne} in die Lösungen zu schauen. Schreiben
Sie Ihre Antworten in vollständigen Sätzen --- die IHK-Prüfung verlangt formulierte
Begründungen, nicht bloße Stichworte. Beachten Sie die \emph{Operatoren} (nennen,
erläutern, beschreiben, begründen, beurteilen) --- sie geben den geforderten
Antwort-Tiefegrad vor. Erst danach öffnen Sie den Lösungsteil ab Seite~\pageref{loesungen}
und vergleichen.
\end{merksatz}

\bigskip

\textbf{Kontext:} Sie sind Auszubildende bei der \textbf{JIKU IT-Solutions GmbH} in Hamburg.
Heute halten die Teams ihre vorbereiteten Betriebspräsentationen. Im Anschluss geben Sie
sich gegenseitig Feedback und wiederholen gemeinsam den gesamten Stoff des Lernfeldes~1.

\bigskip

\noindent
\begin{tabular}{|l|l|l|l|l|l|l|l|}
\hline
\textbf{Aufgabe} & 1 & 2 & 3 & 4 & 5 & 6 & \textbf{Gesamt}\\
\hline
\textbf{Punkte} & /10 & /8 & /8 & /12 & /12 & /10 & /60\\
\hline
\end{tabular}

\tableofcontents
\vspace{1em}
\hrule

\section{Aufgaben}

\subsection*{Aufgabe~1 --- Beobachtung und Feedback geben (10 Punkte)}

\begin{enumerate}[label=(\alph*)]
  \item \textbf{Beobachten} Sie während der Präsentation einer Mitschülergruppe die Vortragstechnik.
  Wählen Sie \textbf{zwei Beobachtungsschwerpunkte} aus den folgenden Bereichen:
  Blickkontakt, Lautstärke und Artikulation, Körperhaltung, Pausen und Tempo,
  Umgang mit Lampenfieber. \textbf{Notieren} Sie zu jedem Schwerpunkt \textbf{mindestens zwei
  konkrete Beobachtungen}. Formulieren Sie nur Beschreibungen --- keine Bewertungen. \hfill (6~P.)

  \bigskip

  \noindent
  \begin{tabular}{|p{3.5cm}|p{9.5cm}|}
  \hline
  \textbf{Schwerpunkt} & \textbf{Konkrete Beobachtungen}\\
  \hline
  & \\[3.5cm]
  \hline
  & \\[3.5cm]
  \hline
  \end{tabular}

  \bigskip

  \item \textbf{Formulieren} Sie auf Basis Ihrer Beobachtungen ein konstruktives Feedback
  (mindestens vier Sätze). Beachten Sie: konkret, beschreibend, als Ich-Botschaft,
  mit Positivem beginnen, nur auf Veränderbares beziehen. \hfill (4~P.)

  \vspace{7cm}
\end{enumerate}

\subsection*{Aufgabe~2 --- Reflexion der Teamarbeit (8 Punkte)}

\begin{enumerate}[label=(\alph*)]
  \item Was ist Ihrer Gruppe bei der Vorbereitung und Durchführung besonders gut gelungen?
  \textbf{Nennen} Sie mindestens zwei konkrete Aspekte. \hfill (2~P.)

  \vspace{4cm}

  \item Was lief weniger gut, und welche Ursache \textbf{vermuten} Sie dafür? \hfill (2~P.)

  \vspace{4cm}

  \item Was würden Sie beim nächsten Präsentationsprojekt konkret anders machen?
  \textbf{Formulieren} Sie eine SMART-Maßnahme (was genau, wer, bis wann). \hfill (2~P.)

  \vspace{4cm}

  \item Welche Reflexionsmethode würden Sie für eine strukturierte Teamauswertung einsetzen ---
  Blitzlicht oder Spinnennetz? \textbf{Begründen} Sie Ihre Wahl. \hfill (2~P.)

  \vspace{4cm}
\end{enumerate}

\subsection*{Aufgabe~3 --- Präsentationsbewertung (8 Punkte)}

\begin{enumerate}[label=(\alph*)]
  \item \textbf{Nennen} Sie die vier Bewertungsbereiche eines typischen
  Präsentations-Bewertungsbogens und die jeweilige Gewichtung. \hfill (4~P.)

  \bigskip

  \noindent
  \begin{tabular}{|p{7cm}|p{5.5cm}|}
  \hline
  \textbf{Bewertungsbereich} & \textbf{Gewichtung}\\
  \hline
  & \\[0.8cm]
  \hline
  & \\[0.8cm]
  \hline
  & \\[0.8cm]
  \hline
  & \\[0.8cm]
  \hline
  \end{tabular}

  \bigskip

  \item \textbf{Erläutern} Sie, welchen Mehrwert der Vergleich von Selbst- und Fremdbewertung
  nach einer Präsentation hat. \hfill (2~P.)

  \vspace{4cm}

  \item Warum muss ein Bewertungsbogen \textbf{vor} der Präsentation verteilt werden?
  \textbf{Begründen} Sie kurz. \hfill (2~P.)

  \vspace{3.5cm}
\end{enumerate}

\subsection*{Aufgabe~4 --- Gesamtwiederholung Ausbildungsrecht (12 Punkte)}

\textbf{Ordnen} Sie die folgenden Aussagen dem richtigen Gesetz und Paragrafen zu. \hfill (12~P.)

\bigskip

\noindent
\begin{tabular}{|p{9.5cm}|p{3.5cm}|}
\hline
\textbf{Aussage} & \textbf{Gesetz / §}\\
\hline
(a) Die Probezeit beträgt mindestens 1 und höchstens 4 Monate. & \\[0.7cm]
\hline
(b) Jugendliche dürfen täglich nicht mehr als 8 Stunden beschäftigt werden. & \\[0.7cm]
\hline
(c) Die Firma ist der Name, unter dem ein Kaufmann seine Geschäfte betreibt. & \\[0.7cm]
\hline
(d) Das Mindest-Stammkapital der GmbH beträgt 25.000~€. & \\[0.7cm]
\hline
(e) Der wesentliche Inhalt des Ausbildungsvertrags ist spätestens vor Ausbildungsbeginn in Textform niederzulegen (seit 1.\,8.\,2024). & \\[0.7cm]
\hline
(f) Die JAV ist die gewählte Interessenvertretung der Jugendlichen und Auszubildenden. & \\[0.7cm]
\hline
\end{tabular}

\subsection*{Aufgabe~5 --- Wertschöpfungsrechnung und Marktanalyse (12 Punkte)}

\textit{JIKU Hamburg erzielt im laufenden Jahr Nettoumsatzerlöse von 4,8~Mio.~€.
Vorleistungen (eingekaufte Hardware und Fremdleistungen) betragen 1,9~Mio.~€,
Abschreibungen 0,3~Mio.~€. Am Hamburger Standort arbeiten 52~Mitarbeiter.}

\begin{enumerate}[label=(\alph*)]
  \item \textbf{Berechnen} Sie die Wertschöpfung. \hfill (2~P.)

  \vspace{3.5cm}

  \item \textbf{Berechnen} Sie die Wertschöpfung je Mitarbeiter. \hfill (2~P.)

  \vspace{3cm}

  \item \textbf{Erläutern} Sie, an welche vier Gruppen die Wertschöpfung verteilt wird. \hfill (4~P.)

  \vspace{5cm}

  \item Der Cloud-Infrastrukturmarkt wird von AWS, Microsoft Azure und Google Cloud beherrscht.
  \textbf{Benennen} Sie die Marktform und \textbf{beschreiben} Sie ein typisches Merkmal
  dieses Markttyps. \hfill (4~P.)

  \vspace{5cm}
\end{enumerate}

\subsection*{Aufgabe~6 --- Fallaufgabe Ausbildungsrecht bei JIKU (10 Punkte)}

\textit{Lena (17 Jahre) beginnt ihre Ausbildung zur Fachinformatikerin bei JIKU Hamburg.
Ihr Ausbilder plant: 9,5 Stunden täglich (8:00--17:30 Uhr), 24 Urlaubstage pro Jahr,
Einreichung des Ausbildungsvertrags bei der IHK erst nach zwei Wochen.}

\begin{enumerate}[label=(\alph*)]
  \item Welche der drei geplanten Regelungen verstoßen gegen geltendes Recht?
  \textbf{Nennen} Sie jeweils das Gesetz und den Paragrafen sowie die korrekte Regelung. \hfill (6~P.)

  \vspace{8cm}

  \item Was versteht man unter der \textbf{gestreckten Abschlussprüfung} der IT-Berufe,
  und wie wird sie gewichtet? \textbf{Erläutern} Sie. \hfill (4~P.)

  \vspace{6cm}
\end{enumerate}

\newpage
\section*{Lösungshinweise für Lehrende}\label{loesungen}

\subsection*{Aufgabe~1}

\textbf{(a)} Beobachtungen sollen konkret und beschreibend sein --- keine Urteile.
Beispiel Blickkontakt: „Die Vortragenden haben während der ersten drei Minuten überwiegend
auf die Folie geschaut" statt „Sie hatten schlechten Blickkontakt."
Akzeptiert werden alle Beobachtungen, die eine wahrnehmbare Handlung beschreiben.

\textbf{(b)} Musterstruktur: (1) Mit Positivem beginnen: Ich-Botschaft über eine Stärke.
(2) Zweite Stärke benennen. (3) Konkrete Beobachtung zu einem Verbesserungsbereich:
„Ich habe bemerkt, dass \ldots" (4) Konkreter Verbesserungsvorschlag.
Alle vier Feedbacksätze müssen als Ich-Botschaften formuliert sein.

\subsection*{Aufgabe~2}

\textbf{(a)} Beispiele: Gute Aufgabenverteilung, klarer Aufbau, Probevortrag durchgeführt.
\textbf{(b)} Beispiele: Zeitplanung unterschätzt, zu viel Text auf Folien, fehlende Abstimmung.
\textbf{(c)} SMART-Beispiel: „Wir halten bis zur nächsten Präsentation in drei Wochen einen
Probevortrag vor der Klasse --- Verantwortlich: alle Gruppenmitglieder."
\textbf{(d)} Blitzlicht: schnell, für kurze Zeitfenster. Spinnennetz: differenzierter,
für Mehrkriterien-Bewertung auf Skala. Beide Antworten akzeptabel bei überzeugender Begründung.

\subsection*{Aufgabe~3}

\textbf{(a)} Fachlicher Inhalt 40~\%, Struktur 20~\%, Medien und Folien 20~\%,
Vortragsstil 20~\%.
\textbf{(b)} Blinde Flecken sichtbar machen: eigene Stärken werden manchmal unterschätzt,
Schwächen nicht erkannt. Nur durch Gegenüberstellung wird Reflexion möglich.
\textbf{(c)} Alle Beurteilenden nutzen dieselben Maßstäbe; die Vortragenden können
sich gezielt vorbereiten; Willkür in der Beurteilung wird verhindert.

\subsection*{Aufgabe~4}

(a) § 20 BBiG --- Probezeit.
(b) § 8 JArbSchG --- tägliche Arbeitszeit.
(c) § 17 HGB --- Firma.
(d) § 5 GmbHG --- Mindeststammkapital.
(e) § 11 BBiG i.\,d.\,F.\ seit 1.\,8.\,2024 (BVaDiG) --- Vertragsniederschrift in Textform; die Kündigung bleibt dagegen schriftformpflichtig (§ 22 Abs.\,3 BBiG).
(f) § 70 BetrVG --- JAV.

\subsection*{Aufgabe~5}

\textbf{(a)} Wertschöpfung = 4.800.000~€ $-$ 1.900.000~€ $-$ 300.000~€ = \textbf{2.600.000~€}.
\textbf{(b)} Wertschöpfung je MA = 2.600.000~€ $\div$ 52 = \textbf{50.000~€}.
\textbf{(c)} Verteilung an: Mitarbeiter (Löhne/Gehälter), Fremdkapitalgeber (Zinsen),
Staat (Steuern und Abgaben), Eigenkapitalgeber (Gewinn/Dividende).
\textbf{(d)} Angebotsoligopol: wenige Anbieter (AWS, Azure, Google Cloud) stehen vielen
Nachfragern gegenüber. Typisches Merkmal: gegenseitige Beobachtung der Wettbewerber,
Preisreaktionen aufeinander abgestimmt, starke Marktmacht der wenigen Anbieter.

\subsection*{Aufgabe~6}

\textbf{(a)} Drei Verstöße:
(1) 9,5~h täglich: Verstoß gegen § 8~JArbSchG --- max. 8~h täglich für Personen unter 18~J.
(2) 24 Urlaubstage: Verstoß gegen § 19~JArbSchG --- Lena ist 17 (noch nicht 18), daher mind. 25 Werktage.
(3) Verspätete IHK-Einreichung: Verstoß gegen § 11~Abs.~4~BBiG --- unverzüglich einzureichen.

\textbf{(b)} Die gestreckte Abschlussprüfung besteht aus zwei Teilen:
Teil~1 (im 4.~Ausbildungshalbjahr, nach ca. 18~Monaten) zählt 20~\% der Gesamtnote;
Teil~2 (am Ende der Ausbildung) zählt 80~\% der Gesamtnote.
Teil~2 umfasst bei Fachinformatikern u.\,a. die Projektarbeit und das Fachgespräch.

\end{document}
